% Generated draft: requires mathematical and source review.
\documentclass[11pt]{article}
\usepackage[margin=25mm]{geometry}
\usepackage{fontspec}
\setmainfont{Latin Modern Roman}
\usepackage{amsmath,amssymb,mathtools,unicode-math}
\setmathfont{Latin Modern Math}
\usepackage{enumitem,needspace,xurl,hyperref}
\setlength{\parindent}{0pt}
\setlength{\parskip}{5pt}
\setlength{\emergencystretch}{3em}
\newcommand{\N}{\mathbb{N}}
\newcommand{\Z}{\mathbb{Z}}
\newcommand{\Q}{\mathbb{Q}}
\newcommand{\R}{\mathbb{R}}
\newcommand{\C}{\mathbb{C}}
\newcommand{\F}{\mathbb{F}}
\newcommand{\A}{\mathbb{A}}
\newcommand{\PP}{\mathbb{P}}
\newcommand{\E}{\mathbb{E}}
\newcommand{\eps}{\varepsilon}
\newcommand{\dd}{\,\mathrm{d}}
\begin{document}
\section*{UnsolvedMath ID: 30005605}
\textbf{Source problem code:} catalogue entry
\subsection{Definitions and mathematical statement}
Let \(S^2=\{x\in\R^3:|x|=1\}\) carry its standard round metric and area measure \(\dd A\), so that \(|S^2|=4\pi\). For \(0<a\leq 2\pi\), let \(C_a\) be a geodesic spherical cap of area \(a\); its geodesic radius is \(\arccos(1-a/(2\pi))\).

For a connected open domain \(\Omega\subset S^2\) whose Neumann Laplacian has discrete spectrum, write its eigenvalues, counted with multiplicity, as
\[
0=\lambda_1^N(\Omega)<\lambda_2^N(\Omega)\leq\lambda_3^N(\Omega)\leq\cdots.
\]
Here the Neumann Laplacian is the operator associated with the quadratic form \(u\mapsto\int_\Omega |\nabla_{S^2}u|^2\,\dd A\) on \(H^1(\Omega)\). In this setting,
\[
\lambda_2^N(\Omega)=\inf_{\substack{u\in H^1(\Omega)\setminus\{0\}\\ \int_\Omega u\,\dd A=0}}
\frac{\int_\Omega |\nabla_{S^2}u|^2\,\dd A}{\int_\Omega |u|^2\,\dd A}.
\]
Connectedness makes the second eigenvalue equal to the first positive eigenvalue. The source does not specify boundary regularity or its precise admissible class of domains; the spectral setting above is an explicit working interpretation requiring review.

\textbf{Proposed inequality recorded by the source.} For every \(a\in(0,2\pi]\) and every admissible domain \(\Omega\subset S^2\) with \(|\Omega|=a\), establish
\[
\lambda_2^N(\Omega)\leq\lambda_2^N(C_a),
\]
without imposing simple connectedness or the source's additional condition that \(\Omega\) lie outside a complementary cap. The source does not define that complementary-cap condition precisely.

\textbf{Unresolved source qualification.} The supplied original discussion reports the inequality for simply connected domains through \(94\%\) of the sphere, with extension toward the full sphere posed separately, and for arbitrary domains of area at most half the sphere subject to the complementary-cap condition. The present target is removal of the latter condition. However, the supplied dated literature assessment also reports counterexamples for multiply connected domains. It does not specify their areas or other hypotheses, so their applicability to the displayed inequality cannot be determined from this record. Accordingly, the displayed assertion is a proposed target, not an established theorem or an independently verified open conjecture.
\subsection{Short English statement}
The proposed comparison says that, at every prescribed positive area up to half the round sphere, a spherical cap has the largest second Neumann eigenvalue, without a complementary-cap restriction. The source leaves the admissible domain class unclear, and its report of multiply connected counterexamples requires checking before this target can be treated as open or valid.
\subsection{Sources}
\begin{itemize}
\item Supplied local UnsolvedMath registry record, permanent ID \(30005605\), source problem code \texttt{catalogue entry}, titled ``Removing Complementary-Cap Restrictions for Spherical Neumann Eigenvalues,'' in the Oberwolfach Reports Open Problems collection.
\item Supplied report citation: \emph{Geometric Spectral Theory (2024)}, Oberwolfach Reports, DOI \texttt{10.4171/owr/2023/36}, pp.~2017--2106. \url{https://doi.org/10.4171/owr/2023/36}. The original passage is attributed to Richard S. Laugesen under ``Neumann and Robin eigenvalues on curved surfaces---open problems.'' The registry gives proposed year 2023 and does not separately identify a proposer.
\item Supplied literature reference: ``Oberwolfach Report 36.'' \url{https://doi.org/10.4171/owr/2023/36}.
\item Supplied literature reference: ``Langford and Laugesen primary article.'' \url{https://arxiv.org/abs/2305.11090}.
\item Supplied literature reference: ``Bucur et al. counterexample.'' \url{https://arxiv.org/abs/2503.15385}.
\item Literature assessment dated exactly \texttt{2026-08-22}: supplied dated curation, not independently verified. It labels the problem partially solved, reports cap maximality for simply connected domains through \(94\%\) of the sphere, describes extension beyond that range as open, and reports rigorous counterexamples only for multiply connected domains. The supplied background explicitly requests independent verification of these time-sensitive claims.
\item The registry labels the statement \texttt{corrected\_verified} and says it was reconstructed and checked against the official Oberwolfach report. These are supplied verification claims; no independent source consultation or mathematical verification is asserted here.
\end{itemize}
\end{document}
